% !TEX program = xelatex

\documentclass[12pt,a4paper]{article}

\usepackage{fontspec}
\usepackage{polyglossia}
\setmainlanguage{russian}
\setotherlanguage{english}
\setmainfont{DejaVu Serif}
\newfontfamily\cyrillicfonttt{DejaVu Sans Mono}
\usepackage{amsmath,amssymb,amsfonts,amsthm,mathtools}
\usepackage{geometry}
\usepackage{booktabs}
\usepackage{tabularx}
\usepackage{titlesec}
\usepackage{cite}
\usepackage{microtype}
\usepackage{xcolor}
\usepackage{tcolorbox}
\usepackage{hyperref}

\microtypesetup{expansion=false}
\setlength{\emergencystretch}{3em}

\titleformat{\section}{\large\bfseries}{\thesection}{1em}{}
\titleformat{\subsection}{\normalsize\bfseries}{\thesubsection}{1em}{}
\titleformat{\subsubsection}{\small\bfseries}{\thesubsubsection}{1em}{}

\geometry{
    left=25mm,
    right=20mm,
    top=25mm,
    bottom=25mm
}

\hypersetup{
    colorlinks=true,
    linkcolor=blue,
    citecolor=blue,
    urlcolor=blue
}

% Стилизованные блоки для справочника
\newtcolorbox{theorybox}[1]{
    colback=gray!6!white,
    colframe=gray!75!black,
    fonttitle=\bfseries\small,
    title=#1,
    arc=1.5mm,
    boxrule=0.8pt,
    top=2mm, bottom=2mm, left=3mm, right=3mm
}

\newtcolorbox{warningbox}[1]{
    colback=red!4!white,
    colframe=red!70!black,
    fonttitle=\bfseries\small,
    title=#1,
    arc=1.5mm,
    boxrule=0.8pt,
    top=2mm, bottom=2mm, left=3mm, right=3mm
}

\begin{document}

\section{Симметричность оператора Штурма--Лиувилля и простота спектра: доказательства и нюансы}

Рассматривается классический дифференциальный оператор второго порядка:
\begin{equation}
    \mathcal{L}y \equiv -\frac{d}{dx}\left(p(x)\frac{dy}{dx}\right) + q(x)y, \quad x \in (a, b),
    \label{eq:SL_operator}
\end{equation}
где $p(x) \in C^1[a, b]$, $p(x) > 0$, $q(x) \in C[a, b]$ — вещественные функции. 

Область определения оператора $\mathcal{D}(\mathcal{L})$ состоит из функций $y \in C^2[a, b]$, удовлетворяющих однородным распадающимся краевым условиям:
\begin{align}
    \Gamma_a(y) &\equiv \alpha_1 y(a) + \alpha_2 y'(a) = 0, \quad (\alpha_1^2 + \alpha_2^2 \neq 0), \label{eq:BC_left} \\
    \Gamma_b(y) &\equiv \beta_1 y(b) + \beta_2 y'(b) = 0, \quad (\beta_1^2 + \beta_2^2 \neq 0). \label{eq:BC_right}
\end{align}

% --------------------------------------------------------------------------
\subsection{Доказательство симметричности оператора}

Оператор $\mathcal{L}$ называется \textit{формально самосопряженным (симметричным)}, если для любых двух функций $u, v \in \mathcal{D}(\mathcal{L})$ выполняется равенство скалярных произведений в $L_2(a, b)$:
\begin{equation}
    \langle \mathcal{L}u, v \rangle = \langle u, \mathcal{L}v \rangle \quad \iff \quad \int_a^b (\mathcal{L}u \cdot v - u \cdot \mathcal{L}v)\,dx = 0.
\end{equation}

\begin{proof}
Вычислим разность подынтегральных выражений, дважды интегрируя по частям (вторая формула Грина / тождество Лагранжа):
\begin{align}
    v \mathcal{L}u - u \mathcal{L}v &= -v (p u')' + u (p v')' \nonumber \\
    &= \frac{d}{dx}\left[ p(x) \left( u(x)v'(x) - u'(x)v(x) \right) \right] = \frac{d}{dx}\left[ p(x) W(u, v)(x) \right],
\end{align}
где $W(u, v)(x) = u(x)v'(x) - u'(x)v(x)$ — определитель Вронского (вронскиан) пары функций $u$ и $v$.

Интегрируя полученное выражение от $a$ до $b$, получаем:
\begin{equation}
    \int_a^b (v \mathcal{L}u - u \mathcal{L}v)\,dx = \left. p(x) W(u, v)(x) \right|_a^b = p(b)W(u, v)(b) - p(a)W(u, v)(a).
    \label{eq:Green_identity}
\end{equation}

Покажем, что каждое из слагаемых в правой части \eqref{eq:Green_identity} равно нулю в отдельности. 
Так как $u, v \in \mathcal{D}(\mathcal{L})$, обе функции удовлетворяют условию \eqref{eq:BC_left} на левом конце:
\begin{equation}
    \begin{cases}
        \alpha_1 u(a) + \alpha_2 u'(a) = 0, \\
        \alpha_1 v(a) + \alpha_2 v'(a) = 0.
    \end{cases}
\end{equation}
Это однородная система линейных алгебраических уравнений относительно коэффициентов $(\alpha_1, \alpha_2)$. Поскольку вектор коэффициентов нетривиален ($|\alpha_1| + |\alpha_2| \neq 0$), определитель матрицы системы обязан быть равен нулю:
\begin{equation}
    \begin{vmatrix} u(a) & u'(a) \\ v(a) & v'(a) \end{vmatrix} = u(a)v'(a) - u'(a)v(a) = W(u, v)(a) = 0.
\end{equation}
Абсолютно аналогично, в силу выполнения условия \eqref{eq:BC_right} на правом конце, получаем $W(u, v)(b) = 0$.

Следовательно, внеинтегральные члены обращаются в ноль:
\begin{equation}
    \left. p(x) W(u, v)(x) \right|_a^b = p(b) \cdot 0 - p(a) \cdot 0 = 0,
\end{equation}
откуда следует $\langle \mathcal{L}u, v \rangle = \langle u, \mathcal{L}v \rangle$. Оператор $\mathcal{L}$ симметричен.
\end{proof}

% --------------------------------------------------------------------------
\subsection{Ловушка восприятия: почему $W(u, v)(x)$ не равен нулю внутри интервала?}
\label{subsec:trap}


Может возникнуть ошибочная мысль: \textit{«Если в доказательстве симметричности оператора $W(u, v)(a) = 0$, то по теореме Лиувилля вронскиан равен нулю везде, а значит, любые две функции зависимы?»}

\begin{warningbox}{Важная ловушка: поведение вронскиана внутри отрезка}
\textbf{Это неверно!} Теорема Лиувилля--Остроградского применима исключительно к функциям, являющимся решениями \textbf{ОДНОГО И ТОГО ЖЕ} линейного однородного уравнения. 

В доказательстве симметричности функции $u(x)$ и $v(x)$ — произвольные функции из области определения $\mathcal{D}(\mathcal{L})$. Они \textbf{не обязаны} быть решениями какого-либо уравнения, поэтому зануление вронскиана на концах не влечет его зануления внутри интервала.

\end{warningbox}

\textbf{Контрпример:} на отрезке $[0, \pi]$ функции $u(x) = \sin(x)$ и $v(x) = \sin(2x)$ удовлетворяют условиям $y(0)=y(\pi)=0$. При этом:
\begin{itemize}
    \item на границах: $W(u, v)(0) = 0$ и $W(u, v)(\pi) = 0$;
    \item в центре ($x = \pi/2$): $W(u, v)(\pi/2) = -2 \neq 0$.
\end{itemize}
Теорема Лиувилля здесь не работает, поскольку $\sin(x)$ и $\sin(2x)$ отвечают \textbf{разным} значениям спектрального параметра ($\lambda = 1$ и $\lambda = 4$).


% --------------------------------------------------------------------------
\subsection{Доказательство однократности (простоты) собственных значений}

\begin{theorybox}{Теорема об однократности (простоте) собственных значений}
Каждому собственному значению $\lambda$ задачи Штурма--Лиувилля с распадающимися краевыми условиями \eqref{eq:BC_left}--\eqref{eq:BC_right} соответствует ровно одна линейно независимая собственная функция (кратность $\lambda$ строго равна 1).
\end{theorybox}

\paragraph{Доказательство.}
Пусть некоторому значению $\lambda$ соответствуют две собственные функции: $y_1(x)$ и $y_2(x)$. 
По определению собственной функции, они обе удовлетворяют \textbf{одному и тому же} дифференциальному уравнению:
\begin{equation}
    -(p(x)y')' + q(x)y = \lambda w(x)y \implies y'' + \frac{p'}{p}y' + \frac{\lambda w - q}{p}y = 0,
    \label{eq:ODE_fixed_lambda}
\end{equation}
а также краевому условию в точке $x = a$:
\begin{equation}
    \begin{cases}
        \alpha_1 y_1(a) + \alpha_2 y'_1(a) = 0, \\
        \alpha_1 y_2(a) + \alpha_2 y'_2(a) = 0.
    \end{cases}
\end{equation}
Из нетривиальности вектора $(\alpha_1, \alpha_2)$ строго следует равенство нулю определителя системы:
\begin{equation}
    W(y_1, y_2)(a) = 0.
\end{equation}
Поскольку функции $y_1$ и $y_2$ являются решениями \textbf{одного и того же уравнения} \eqref{eq:ODE_fixed_lambda}, к ним применима формула Лиувилля--Остроградского:
\begin{equation}
    W(y_1, y_2)(x) = W(y_1, y_2)(a) \cdot \exp\left( -\int_a^x \frac{p'(\xi)}{p(\xi)}\,d\xi \right) = 0 \cdot \frac{p(a)}{p(x)} \equiv 0 \quad \forall x \in [a, b].
\end{equation}
Так как вронскиан двух решений линейного однородного ОДУ второго порядка тождественно равен нулю на интервале, эти решения \textbf{линейно зависимы}:
\begin{equation}
    y_2(x) = C \cdot y_1(x), \quad C = \text{const}.
\end{equation}
Следовательно, размерность собственного подпространства равна 1. $\blacksquare$

% --------------------------------------------------------------------------
\subsection{Почему для периодических условий кратность может быть равна двум?}

Рассмотрим краевые условия периодического типа:
\begin{equation}
    y(a) = y(b), \qquad y'(a) = y'(b).
    \label{eq:BC_periodic}
\end{equation}
Условия \eqref{eq:BC_periodic} по-прежнему обеспечивают самосопряженность оператора, так как:
\begin{equation*}
    \left. p(x) W(u, v)(x) \right|_a^b = p(b)W(u, v)(b) - p(a)W(u, v)(a) = 0 \quad (\text{при } p(a)=p(b)).
\end{equation*}

Однако доказательство однократности спектра здесь не работает.
Причина: периодические условия связывают значения на противоположных границах, но не накладывают никаких локальных ограничений на вектор $(y(a), y'(a))$. В точке $x = a$ решение сохраняет обе степени свободы. Следовательно, для двух собственных функций $y_1, y_2$ их вронскиан в точке $a$ не обязан быть равен нулю ($W(y_1, y_2)(a) \neq 0$).

\paragraph{Пример.} Для задачи $-y'' = \lambda y$ на $[-\pi, \pi]$ с периодическими условиями каждому собственному значению $\lambda_n = n^2$ ($n \ge 1$) соответствуют две линейно независимые собственные функции:
\begin{equation*}
    y_{n, 1}(x) = \cos(nx) \quad \text{и} \quad y_{n, 2}(x) = \sin(nx).
\end{equation*}
Их вронскиан $W(\cos nx, \sin nx) = n \neq 0$ везде, включая концы отрезка. Кратность собственного значения равна 2.

% --------------------------------------------------------------------------
\subsection{Сводная классификация краевых условий и их последствий}

В таблице \ref{tab:BC_classification} приведена систематизация различных вариантов постановки краевых условий для дифференциального оператора 2-го порядка на отрезке $[a, b]$.

\begin{table}[h!]
\small
\centering
\caption{Влияние выбора граничных условий на постановку задачи}
\label{tab:BC_classification}
\vspace{0.5em}
\begin{tabularx}{\textwidth}{p{3.5cm} p{3.2cm} p{2.8cm} X}
\toprule
\textbf{Тип условий} & \textbf{Пример записи} & \textbf{Статус задачи} & \textbf{Следствия для спектра / решений} \\
\midrule
\textbf{2 распадающихся однородных} (Штурм--Лиувилль) & $\alpha_1 y(a) + \alpha_2 y'(a) = 0$ $\beta_1 y(b) + \beta_2 y'(b) = 0$ & Классическая спектральная задача & $\mathcal{L}$ самосопряжен. $\lambda \in \mathbb{R}$. Собственные функции ортогональны. Кратность $\lambda$ строго равна 1. \\
\midrule
\textbf{2 нераспадающихся однородных} (Периодические) & $y(a) = y(b)$ $y'(a) = y'(b)$ & Спектральная задача & $\mathcal{L}$ самосопряжен. $\lambda \in \mathbb{R}$. Возможна двукратная кратность собственных значений. \\
\midrule
\textbf{Неоднородные условия} & $y(a) = A$ $y(b) = B$ ($|A|+|B| \neq 0$) & Не является спектральной задачей & Обычная краевая задача. Решение $y \equiv 0$ невозможно. Требуется сведение к однородным замене $y = u + g$. \\
\midrule
\textbf{3 или 4 условия} (Переопределенная система) & $y(a) = 0, y'(a) = 0,$ $y(b) = 0$ & Некорректная задача (переопределена) & Пространство решений ОДУ 2-го порядка двумерно. Задача Коши $y(a)=y'(a)=0$ дает только $y(x) \equiv 0$. Спектр пуст. \\
\midrule
\textbf{1 условие} (Недоопределенная система) & $y(a) = 0$ (на правом конце условий нет) & Некорректная задача (недоопределена) & Не хватает условий для фиксации оператора; спектр не является дискретным; отсутствует единственность решений. \\
\bottomrule
\end{tabularx}
\end{table}

\section{Почему краевые условия в задаче Штурма--Лиувилля обязаны быть однородными?}

Классическая регулярная задача Штурма--Лиувилля формулируется как спектральная задача (задача на собственные значения) для дифференциального оператора второго порядка:
\begin{equation}
    \mathcal{L}y \equiv -\frac{d}{dx}\left(p(x)\frac{dy}{dx}\right) + q(x)y = \lambda w(x)y, \quad x \in (a, b),
    \label{eq:SL_equation}
\end{equation}
где $p(x) > 0$, $w(x) > 0$, $p'(x), q(x), w(x)$ непрерывны на $[a, b]$, с краевыми условиями:
\begin{align}
    \alpha_1 y(a) + \alpha_2 y'(a) &= 0, \quad |\alpha_1| + |\alpha_2| \neq 0, \label{eq:BC_a} \\
    \beta_1 y(b) + \beta_2 y'(b) &= 0, \quad |\beta_1| + |\beta_2| \neq 0. \label{eq:BC_b}
\end{align}
Требование равенства нулю правых частей в условиях \eqref{eq:BC_a}--\eqref{eq:BC_b} (однородность) является фундаментальным и обусловлено следующими причинами:

\subsection*{1. Определение спектральной задачи}
По определению, число $\lambda$ называется \textit{собственным значением}, если существует \textbf{нетривиальное} ($y(x) \not\equiv 0$) решение соответствующей задачи. 

При любом значении параметра $\lambda$ функция $y(x) \equiv 0$ тождественно удовлетворяет уравнению \eqref{eq:SL_equation}. Чтобы она являлась решением всей краевой задачи, она обязана удовлетворять и краевым условиям, что возможно \textbf{только} при нулевых правых частях:
\begin{equation*}
    \alpha_1 \cdot 0 + \alpha_2 \cdot 0 = 0.
\end{equation*}
Если задать неоднородное условие (например, $y(a) = 1$), то тривиальное решение исключается, и теряется само определение собственного значения: задача перестает быть спектральной и становится обычной краевой задачей.

\subsection*{2. Линейность области определения оператора}
Для корректного определения оператора $\mathcal{L}$ в гильбертовом пространстве $L_2(a, b; w(x))$ его область определения $\mathcal{D}(\mathcal{L})$ обязана быть \textbf{линейным подпространством}. Это означает, что для любых функций $y_1, y_2 \in \mathcal{D}(\mathcal{L})$ и любых констант $c_1, c_2 \in \mathbb{C}$ их линейная комбинация также должна принадлежать $\mathcal{D}(\mathcal{L})$:
\begin{equation*}
    (c_1 y_1 + c_2 y_2) \in \mathcal{D}(\mathcal{L}).
\end{equation*}
Если краевое условие неоднородно, например, $y(a) = C \neq 0$, то для суммы решений получаем:
\begin{equation*}
    (y_1 + y_2)(a) = y_1(a) + y_2(a) = 2C \neq C.
\end{equation*}
Множество функций, удовлетворяющих неоднородным условиям, является не векторным подпространством, а лишь \textit{аффинным многообразием}. В таком множестве нельзя определить базис из собственных функций.

\subsection*{3. Самосопряженность оператора и зануление внеинтегральных членов}
Ключевое свойство задачи Штурма--Лиувилля — вещественность спектра и ортогональность собственных функций — гарантируется \textbf{самосопряженностью (эрмитовостью)} оператора $\mathcal{L}$.

Рассмотрим формулу Грина (второе тождество Лагранжа) для оператора $\mathcal{L}$ на гладких функциях $u(x)$ и $v(x)$:
\begin{equation}
    \int_a^b \left( u \mathcal{L}v - v \mathcal{L}u \right) dx = \left[ p(x) \left( u(x)v'(x) - u'(x)v(x) \right) \right]_a^b.
    \label{eq:Green_formula}
\end{equation}
Оператор $\mathcal{L}$ будет симметричным относительно стандартного скалярного произведения $\langle u, v \rangle = \int_a^b u(x)v(x)\,dx$ тогда и только тогда, когда внеинтегральный член в \eqref{eq:Green_formula} обращается в ноль:
\begin{equation}
    \left. p(x) W(u, v)(x) \right|_a^b = 0,
    \label{eq:Wronskian_zero}
\end{equation}
где $W(u, v)(x) = u v' - u' v$ — вронскиан функций $u$ и $v$.

Если обе функции $u(x)$ и $v(x)$ удовлетворяют однородному условию \eqref{eq:BC_a} на левом конце:
\begin{equation*}
    \begin{cases}
        \alpha_1 u(a) + \alpha_2 u'(a) = 0, \\
        \alpha_1 v(a) + \alpha_2 v'(a) = 0,
    \end{cases}
\end{equation*}
то эта система имеет нетривиальное решение $(\alpha_1, \alpha_2) \neq (0, 0)$ только при условии равенства нулю её определителя:
\begin{equation*}
    \begin{vmatrix} u(a) & u'(a) \\ v(a) & v'(a) \end{vmatrix} = u(a)v'(a) - u'(a)v(a) = W(u, v)(a) = 0.
\end{equation*}
Аналогично, на правом конце $W(u, v)(b) = 0$. 

При нарушении однородности краевых условий (если в правых частях стоят ненулевые константы) внеинтегральный член \eqref{eq:Wronskian_zero} не обращается в ноль. Это приводит к разрушению симметричности оператора, и, как следствие:
\begin{itemize}
    \item собственные значения $\lambda$ перестают быть гарантированно вещественными;
    \item собственные функции, соответствующие разным собственным значениям, теряют взаимную ортогональность с весом $w(x)$;
    \item теряется теорема разложения (теорема Стеклова) функций в сходящийся ряд Фурье по собственным функциям задачи.
\end{itemize}

\paragraph{Замечание о неоднородных задачах.}
Если физическая постановка содержит неоднородные краевые условия (например, $y(a) = A$, $y(b) = B$), то соответствующая задача является \textit{краевой задачей математической физики}, но \textbf{не задачей Штурма--Лиувилля}. Для применения спектрального аппарата такую задачу всегда предварительно сводят к задаче с однородными условиями с помощью замены переменной $y(x) = u(x) + g(x)$, где $g(x)$ — любая фиксированная функция, удовлетворяющая неоднородным условиям ($g(a)=A, g(b)=B$), а для искомой функции $u(x)$ условия становятся строго однородными: $u(a)=0, u(b)=0$.

\end{document}